\begin{theorem}[Альтернатива Фредгольма]
% см. в книге Колмогорова-Фомина альтернативную формулировку
    \setcounter{equation}{0}
    Исследуем разрешимость уравнений
    \begin{equation}
    A_{\lambda} x = y
    \end{equation}
    \begin{equation}
    A_{\lambda} z = 0
    \end{equation}
    
    Пусть $\lambda \neq 0$. Две альтернативы
    \begin{itemize}
        \item $\lambda$~--- собственное значение оператора $A$.
        Откуда следует, что $\ker A_\lambda \neq 0$, значит образ - это не все пространство. Уравнение (1) имеет решение (не единственное) $\lra y \perp$ всем решениям уравнения (2)
        \item $\lambda$~--- не собственное значение.
        
        По лемме \nameref{fredholm:lemma1} $\lambda \in \rho(A)$. А значит, 
        уравнение (1) имеет единственное решение при любой правой части.
    \end{itemize}

\end{theorem}

% тут картинка от Поли про взаимосвязи теорем

% Коновалов приехал в страну компактных операторов и ничего там не нашел

\begin{theorem}[12.5 (Гильберт-Шмидт)]\label{12.5}
    $H(\C)$~--- сепарабельное гильбертово пространство, $A: H \to H$~--- компактный самосопряжённый оператор. Тогда в $H$ существует ортонормированный базис, состоящий из собственных векторов оператора $A$.
\end{theorem}

\begin{lemma}[4]\label{hilbert-shmidt:lemma4}
        Пусть $H(\C)$~--- гильбертово пространство, $A$~--- компактный самосопряжённый оператор, $A \neq 0$.
        Тогда у $A$ существует собственное значение, отличное от нуля.
\end{lemma}

\begin{proof}
    По теореме \nameref{11.5} $r(A)=max(|m_-, m_+|)= \|A\| \neq 0$ (так как оператор ненулевой).
    Так как $A \neq 0$, то хотя бы одно из этих чисел ($m_+, m_-$) не равно нулю. 
    То, которое не ноль, по лемме \nameref{fredholm:lemma1} это $m_* \neq 0$~--- собственное значение.
\end{proof}

\begin{proof}[доказательство теоремы Гильберта-Шмидта]
    Будем использовать теоремы \nameref{11.5}, \nameref{12.3} и лемму \nameref{hilbert-shmidt:lemma4}

    % яблоко хочет чтобы его съели
    
    Вспоминаем, что $A$~--- компактный оператор. По теореме \nameref{12.3} внешность любой окрестности нуля содержит конечное число собственных значений.
    Упорядочим эти собственные значения по убыванию модулей $|\lambda_1| \geqslant |\lambda_2| \geq \ldots$. Этим собственным значениям соответствуют собственные векторы $e_1, e_2, \ldots$. Теперь, уменьшая радиус окрестности нуля, упорядочим все ненулевые собственные значения.

    $\{ e_n \}, \lambda \neq 0$
    $\ker A \subset H$~--- сепарабельное, значит в нем можно выбрать ортонормированный базис $\{ E_0 \}$. Обозначим $\{e\} = \{e_n\} \cup \{E_0\}$. Все элементы $\{e\}$ ортогональны (по теореме \nameref{11.1}, так как собственные векторы соответствующие разным собственным значениям ортогональны).

    $\{e\}$~--- ортонормированная система. Как доказать, что она базис?

    Вспоминаем: ортонормированная система в $H$ является базисом $\lra M:=\overline{[\{e\}]}=H$.
    Запишем теорему о проекции: $M \oplus M^\perp = H$. Нужно доказать, что $M^\perp = \{0\}$.

    $M$ инвариантно относительно $A$: если $x \in M$, то $x = \lim \limits_{n} \sum \limits_k \alpha_{n_k} e_k$. Подействовав оператором $A$, получаем $Ax = \lim \limits_{n} \sum \limits_k \alpha_{n_k} A e_k = \lambda_k e_k$.

    По лемме \nameref{fredholm:lemma3} $M^\perp$ инвариантен относительно $A$.
    Рассмотрим оператор $\tilde{A} = \left. A \right|_{M^\perp}$
    ~--- компактный самосопряжённый оператор.

    Рассмотрим два варианта
    \begin{itemize}
        \item $\tilde{A} = 0$.
        Это означает, что $\tilde{A} x = 0 \: \forall x \in M^\perp$. Следовательно, 
        $M^\perp = \ker \tilde{A}$. В $M^\perp$ могут быть только собственные векторы, соответствующие 0, но все собственные векторы остались в $M$. Если $M^\perp \neq 0$, то приходим к противоречию. Это и означает, что $M^\perp = \{0\}$.
        \item $\tilde{A} \neq 0$.
        Тогда, по лемме \nameref{hilbert-shmidt:lemma4} у $\tilde{A}$ есть собственные значение $\lambda \neq 0$ (оно же является собственным значением оператора $A$).
        У $\lambda$ есть соответсвующий собственный вектор, который должен остаться в $M$, а значит $M^\perp = \{0\}$. % TODO: какая-то фигня написана, поправить
    \end{itemize}
\end{proof}

% Если вы видите только кенгуру не выше 3 метров -- значит вы в компактной Австралии

\begin{example}
    Задача на применение теоремы \nameref{12.5}
    Найти норму оператора Вольтерра 

    \[
    (12-18) (V f)(x) = \int_{t = 0}^x f(t) dt \text{ в } L^2[0, 1]
    \]

\end{example}

Этот подход хорошо соотносится с поиском нормы оператора в $\R_2^n$ (который действует домножением на матрицу $(a_{i_j})$.
\[
\|A\|^2 = (Ax, Ax) = (x, A^*Ax), A^*A  \text{~--- самосопряжённый оператор}
\]

$\{e_k\}$~--- ортонормированный базис.

\[
x = \sum_1^n \xi_k e_k = (\sum \xi_i e_i, \sum \xi_i \lambda e_i) = 
\sum \lambda_i \xi_i^2 \leqslant \max_k \lambda_k \cdot \sum \xi_k^2
\]
(неравенство Парсеваля)

\begin{corollary}
    \begin{enumerate}
        \item $x \in H, \; Ax=A(\sum (x, e_n)e_n)=\sum_{\lambda_n \neq 0} \lambda_n (x, e_n) e_n$
        \item Пусть $\lambda \in \rho(A)$. Решим уравнение $(A - \lambda I)x = y$.
        \[
        \sum \lambda_n (x, e_n) e_n - 
        \sum \lambda (x, e_n) e_n = 
        \sum (y, e_n) e_n
        \]
        Осталось приравнять коэффициенты:
        \[
        (x, e_n) = \frac{(y, e_n)}{\lambda_n - \lambda}
        \]
        \[ R_\lambda = \sum \frac{1}{\lambda_n - \lambda} P_{e_n}\]
        % жующими движениями пробираемся к новым методам
        % вспоминаем то, чего не было
        \item $\lambda = \lambda_k$

        Если $n \neq k$, то формула~--- $(x, e_n) = 
        \frac{(y, e_n)}{\lambda_n - \lambda}$.
        \item Вид спектра компактного самосопряжённого оператора в сепарабельном гильбертовом пространстве

        \begin{enumerate}
            \item[1] $\{\lambda_n\}$~--- бесконечно много
            \item[1.a] $0$~--- собственное значение (может быть любой кратности): $\sigma(A)=\sigma_p$
            \item[1.b] $0$~--- не собственное значение: $\sigma(A)=\sigma_p \cup \{0\}$
            Ноль принадлежит непрерывному спектру
            \item[2] $\{\lambda_n\}$~--- конечно.
            Тогда 0~--- собственное значение бесконечной кратности, $\sigma = \sigma_p$.
        \end{enumerate}
    \end{enumerate}
\end{corollary}

\begin{example}  % TODO доделать
    \begin{enumerate}
        \item[1a] Чередуем 0 и дроби на диагонали
        $$
        \begin{pmatrix}
            0 \\
            0 & 1 \\
            0 & 0 & 0 \\
            0 & 0 & 0 & \frac{1}{2} \\
            \vdots
        \end{pmatrix}
        $$
        \item[1b] На диагонали только дроби
        $$
        \begin{pmatrix}
            1 \\
            0 & \frac{1}{2} \\
            \vdots
        \end{pmatrix}
        $$
        \item[2] Конечное число дробей, потом все нули
    \end{enumerate}
\end{example}

\begin{lemmanote}
    $C[a,b] \subset L_2[a,b]$. $\mathcal{K}=\int_a^b K(x,t)f(t)dt - \lambda f(x)=g(x)$. Если $K \in L_2([a,b]^2) \implies \mathcal{K}$~--- компактный оператор
    Если $K(x, t) = \overline{K(t, x)}$, то это симметричный оператор (мы это обсуждали).
    Тогда к такому оператору применима вся теория для компактных операторов (см §14).
\end{lemmanote}

\begin{example}
    Решаем уравнение для колебания струны
    $$
    \begin{cases}
        u_{tt} = a^2 u_{xx}\\
        \left. u \right|_{t = 0} = u_0(x),
        \left. u_{t = 0} \right|_{t = 0} = u_1(x)
    \end{cases}
    $$
    % дописать философию
\end{example}

\textbf{Операторы Гильберта-Шмидта (см. §14):}
$H = L_2[a, b]$

\[ (
Af)(x) = \int_{t=a}^b K(x,t)f(t)dt, \; K \in L_2([a,b]^2) \implies A \text{~--- компактный оператор} 
\]

% Дан $A$ компактный самосопряжённый опрератор в бесконечномерном гильбертовом пространстве. Входит ли ноль в спектр оператора? (недописано)









